% Lösungen Einheit 1 von 4 – Mittlere Änderungsrate und Sekante (zusammenbau v0.1)
\einheitenkopf{Einheit 1 von 4 · Mittlere Änderungsrate und Sekante}

\erg{10}{a) Zeitraum – Differenzenquotient, Sekante: „in den ersten fünf Stunden“ nennt ein Intervall \quad b) $\frac{300 - 120}{3} = 60$ Besucher je Stunde \quad c) Zeitraum $1{,}5$ h; $\frac{1\,530 - 480}{1{,}5} = 700$ Anmeldungen je Stunde \quad d) $h(5) \approx 63{,}36$, $h(0) = 40$; $\frac{h(5) - h(0)}{5} \approx 4{,}67$ Prozentpunkte je Tag \quad e) $h(5) = 42{,}5$, $h(2) = 50$; $\frac{1}{3} \cdot (42{,}5 - 50) = -2{,}5$: Zwischen der zweiten und der fünften Stunde sinkt der Wasserstand im Mittel um $2{,}5$ cm je Stunde (mittlere Änderungsrate, nicht die momentane). \quad f) $f(1) \approx 2{,}011$, $f(5) \approx 0{,}947$; $m = \frac{f(5) - f(1)}{4} \approx -0{,}266$; $y = -0{,}266 \cdot (x - 1) + 2{,}011 \approx -0{,}266x + 2{,}277$ \quad g) $f$: $\frac{f(5) - f(0)}{5} \approx -0{,}211$; $h$: $\frac{0{,}6 - 1{,}8}{5} = -0{,}24$; der Betrag der mittleren Steigung ist bei $h$ größer ($0{,}24 > 0{,}211$) \quad h) Länge $4 \cdot 3{,}2 + 2\pi \cdot 0{,}5 \approx 15{,}94$ m; Zeit $\frac{15{,}94}{4} \approx 3{,}99$ min \quad i) $\frac{k(c) - k(0)}{c} = 0{,}1c^2 - 1{,}5c = -5$, also $c^2 - 15c + 50 = 0$ und $c = 5$ oder $c = 10$; kleinere Lösung $c = 5$: Die linke Seite ist die mittlere Abbaurate über $[0; c]$; in den ersten fünf Jahren nimmt der Bestand im Mittel um $5$ Mio. Tonnen je Jahr ab. \quad j) Die Sekante von $(0 | 0)$ nach $(3 | 2{,}7)$ verläuft flacher als die $45°$-Gerade $y = x$ – das Kriterium ist erfüllt. \quad k) Falsch: $K(2) - K(1) = 10$, aber $K(3) - K(2) = 4$ – die Kosten des zusätzlichen Kubikmeters nehmen zunächst ab, obwohl $K$ steigt.}
\erg{11}{a) Fehler: Lea gibt die Differenz an, nicht die Rate; die Änderung muss durch die Länge des Zeitraums geteilt werden. Richtig: $\frac{600}{4} = 150$ Besucher je Stunde \quad b) Die Steigung einer Geraden ist Höhenunterschied durch waagerechten Unterschied; für die Sekante sind das $f(b) - f(a)$ und $b - a$ – genau der Differenzenquotient. \quad c) $\frac{380 - 800}{6} = -70$ Liter je Stunde – der Verbrauch liegt im Mittel bei $70$ Litern je Stunde, die Vorgabe ist nicht eingehalten. \quad d) Punkte $(1 | 3)$ und $(3 | 8)$; Sekantensteigung $\frac{8 - 3}{2} = 2{,}5$ cm je Stunde}
